\bookmarksetup{startatroot}
\chapter{Postfață}

Sper că parcurgerea acestui volum a fost și pentru voi o plimbare fascinantă prin lumea algoritmilor și a structurilor de date, după cum și pentru mine a fost o imensă plăcere să îl scriu.

Dar orizonturile sînt mai largi de atît. De aceea, aș vrea să vă las cîteva gînduri de final. De fapt unul singur.

Noi nu facem haine, înghețată, case sau bărci. Toate aceste lucruri sînt bune, dar influența lor asupra lumii în ansamblu este limitată. Noi scriem software; iar software-ul are potențialul să schimbe modul în care lumea gîndește și se comportă, cu efecte pe ani sau pe generații.

Iată cîteva exemple de efecte ale informatizării societății, cu bune și cu rele.

Planeta este mai mică. Comunicarea oriunde pe glob este facilă. Cînd eram eu student în Boston, sunam acasă cu \$2/minut. Astăzi este gratis. Dar ușurința accesului depinde și de cooperarea furnizorilor (ISPs). Furnizorii nu vor ca Internetul să fie considerat o \textit{utilitate publică}, la fel ca apa sau curentul. Ei vor să aibă monopolul asupra infrastructurii, dar să fie liberi să practice orice tarife doresc, cum ar fi tarife diferențiate pentru fiecare aplicație. De exemplu, în 2025 a picat \href{https://en.wikipedia.org/wiki/Net_neutrality}{net neutrality}, încercarea democraților din SUA de a reglementa furnizorii ca pe utilități. Vom vedea care vor fi efectele

Avem standarde și protocoale pentru tot. Vă dați seama ce coșmar ar fi dacă fiecare site ar avea propriul său protocol? Dacă am avea nevoie de programe diferite ca să vizităm Wikipedia sau YouTube? Iar standardele (fie ele HTTP, HTML, XML, PDF, ODF...) există pentru că oamenii au militat pentru asta. Rezultatul este un internet descentralizat și (în general) universal accesibil. Necazul este că această deschidere se năruie în prezent. De exemplu, Facebook a prosperat la începuturi pentru că le-a oferit primilor săi utilizatori interfațarea cu Myspace, o fostă rețea socială. La rîndul său Facebook, acum că este într-o poziție dominantă, face eforturi considerabile ca să nu le permită altor nou-veniți să se integreze cu Facebook. Adobe încearcă de ani buni să monopolizeze formatul PDF, să încarce protocolul cu extensii nelibere care pot fi citite doar cu Adobe Reader (și nu sub GNU/Linux).

Avem calculatoare interoperabile (în general). Dacă descărcăm distribuția oficială a sistemului nostru de operare favorit, o putem instala pe orice nou calculator. Pericolul provine de la treacherous computing (numit impropriu și trusted computing). Acesta promite să ne ferească calculatorul de software neautorizat. Dar subtextul este „neautorizat de producător”. Sistemul TPM (treacherous platform module)  face calculatorul să asculte de producător în detrimentul utilizatorului.

Avem software liber pentru majoritatea nevoilor informatice. Aceasta deși toată lumea se întreabă, de 50 de ani, „de unde cîștigă oamenii ăștia bani?”. Întrebarea este însă prost pusă. Oamenii scriu software și aleg să-l publice sub o licență liberă pentru o mulțime de motive care nu țin de bani. Oamenii fac fapte bune pentru că e în natura umană. Sau pentru că au primit ceva gratuit la rîndul lor și se simt datori. Sau pentru că vor să-și cîștige nemurirea. Sau ca să cîștige inima unei fete / unui băiat. Dar motivul pentru care avem software liber pentru majoritatea nevoilor este că diverși oameni cu spirit de meștereală au dezasamblat drivere binare, au făcut \textit{reverse engineering} la formate și la protocoale și au publicat versiuni libere ale acelorași unelte. Pericolul, astăzi, provine din \href{https://pluralistic.net/tag/drm/}{DRM} (\textit{digital restrictions management}, numit impropriu și \textit{digital rights management}). DRM nu este o unealtă informatică. Este o unealtă legală care incriminează încercarea de a ocoli anumite restricții digitale.

Avem, în era digitală, acces la mai multe opere artistice și tehnice decît orice altă generație. Pericolul vine din durata tot mai lungă a restricțiilor de copyright și de patente. În plus, în cazul software-ului, există pericolul licențelor draconice. Mulți dintre noi le acceptăm pe necitite, dar mesajul lor este dur: ne fac să promitem să fim niște prieteni răi, să nu împărțim acel produs cu nimeni. Un gest firesc acum 40 de ani (să-i împrumuți o carte sau o dischetă \textit{floppy} prietenului tău) a fost îngreunat tehnic și stigmatizat din punct de vedere moral, transformîndu-i în criminali pe cei care știu să fenteze cătușele digitale. Voi sînteți generația digitală, care poate copia terabytes întregi cu o apăsare de taste, și pe care totuși marii deținători de copyright încearcă să vă convingă, prin îndoctrinare zilnică, că copierea este un act rușinos.

Sîntem generația cea mai avansată științific. Pe mine acest lucru mă umple de respect față de lungul lanț al oamenilor de știință care au trăit și au murit pentru ca noi să putem cunoaște universul în măsura în care îl cunoaștem astăzi. Dar, în același timp, Internetul de astăzi constă dintr-un \textit{whitelist} -- o mînă de site-uri cu reputație din care ne putem informa corect, -- și o hazna de site-uri dezinformatoare și antiștiințifice.

\begin{center}
  $\ast$ \quad $\ast$ \quad $\ast$
\end{center}

De ce am enumerat aceste lucruri bune și rele? Pentru că voi sînteți cu toții elevi de excepție și veți avea porți deschise la toate companiile notabile. Vă încurajez ca, oriunde vă veți angaja, să nu renunțați la filtrul moral. Să vă puneți, la fiecare nou proiect, întrebarea: acest feature pe care mi-a cerut compania să-l implementez este un pas în direcția bună, sau un pas în direcția greșită?

Dacă angajatorul vostru vă cere să lucrați la o aplicație care neagă încălzirea globală sau teoria evoluției; sau care identifică persoane care au discutat pe Internet despre anumite cuvinte-cheie ca să-i predea autorităților; sau care verifică dacă piesele de schimb folosite pentru un telefon sînt originale și \href{https://www.wired.com/story/apple-iphone-parts-pairing-right-to-repair/}{refuză să le folosească} în caz contrar; sau care programează un automobil să iasă în stradă și \href{https://news.ycombinator.com/item?id=29623658}{să sune din claxon} imediat ce proprietarul a omis să plătească rata, pentru ca recuperatorul să poată găsi mai ușor automobilul; în toate aceste cazuri sper să luați în calcul să vă căutați de lucru altundeva.

Vă promit că nu veți fi muritori de foame dacă vă veți autoimpune niște standarde de etică. Veți fi, poate, nevoiți să acceptați o ofertă cu 10\% mai mică. Dar o conștiință curată merită deranjul. \emoji{slightly-smiling-face}
